\documentclass[paper=a4,fontsize=11pt]{temp} % KOMA-article class			
\usepackage{eurosym}
\usepackage{float}
\begin{document}
\pagenumbering{alph}
\pagestyle{headings}

\begin{minipage}{.2\linewidth}
   
\includegraphics[width=1\textwidth]{photo}
\end{minipage}      
\begin{minipage}{0.7\linewidth}
   \MyName{ Oguz Umut  SALMAN}
   \sepspace
   \noindent
   
   \hfill umut.salman@lspm.cnrs.fr  
   
   \hfill www.oguzumutsalman.org
   
   \hfill +33(0)783146245   
 
\end{minipage}
%

\NewPart{PROFESSIONAL EXPERIENCE}{}
\noindent


\workEntry{CNRS}{01.2015 – PRESENT}{Chargé de recherche}{Researcher - Centre National Research Scientifique (CNRS) in Mechanics, FRANCE
Lecturer - Université Sorbonne Paris Nord, FRANCE
.} 
{photo}

\sepspace

\workEntry{CNR-IENI }{09.2013 – 09.2015 }{Scientific researcher}{Post-Doctoral Fellow - CNR-IENI Milano, ITALY
Advanced Grant of the European Research Council: Size Effects in Fracture and Plasticity
.} {ieni2}

\sepspace

\workEntry{Bosphorus University}{10.2011 – 10.2012}{Post-doc}{Post-Doctoral Fellow - The Physics Department under the supervision of M. Mungan.}{bos}

\sepspace

\workEntry{Harvard University}{10.2012 – 09.2013}{Post-doc}{Post-Doctoral Fellow - The Aisenberg Biomineralization and Biomimetic Laboratory, Harvard University, USA.}{harvard}

\sepspace

\workEntry{Ecole Polytechnique}{06.2009 – 10.2011}{Post-doc}{Post-Doctoral Fellow - Solid Mechanics Laboratory (LMS), Ecole Polytechnique, FRANCE}{ept}

\NewPart{Education}{}
\noindent



\EducationEntry{Phd. in Physics }{09.2005 – 06.2009}{Université Pierre et Marie Curie, Paris, FRANCE  (High Honour Mention).}{} {sor}

\sepspace

\EducationEntry{MSc. in Materials Science and Nano-Objects }{09.2004 – 09.2005}{Ecole Polytechnique, Palaiseau, FRANCE  (High Honour Mention).} {}{ept}


\sepspace

\EducationEntry{B.S. in Physics  }{09.2000 – 06.2004}{Istanbul Technical University, Physics Department, Istanbul, TURKEY (Honour List).} {}{itu}

\sepspace

\NewPart{Skills \& Interests}{}
\hspace{3mm}
\begin{minipage}[t]{0.33\textwidth} 

\begin{tabular}[t]{ l l }
\flag{IMG/flag/turkish}  & Native Speaker \\
\flag{IMG/flag/fr}  & Fluent \\
\flag{IMG/flag/gb}  & Fluent \\
%\flag{IMG/flag/de}  & Conversational level \\
%\flag{IMG/flag/es}  & Conversational level \\
\end{tabular}

\sepspace

\end{minipage}
%
\begin{minipage}[t]{0.66\textwidth} 


\begin{tabular}[t]{l l}
\textbf{PROGRAMMING}: C, C++, Fortran, Mathematica, Matlab, \\Parallel coding (OpenMP-MPI), ABAQUS ,\\ FreeFem++, GPU Programming\end{tabular}

\end{minipage}
\newpage
\NewPart{Publications}{}
\hspace{3mm}
 O.U. Salman, M Ghidelli, J. Smiri  and I.R. Ionescu, "\textit{Computing the true stress in micro-pillars compression experiments
"}, Submitted to Scripta Materials,  (2023)\\


R. Baggio, O.U. Salman, L. Truskinovsky, "\textit{Nucleation of dislocations by pattern formation"}, European Journal of Mechanics-A/Solids 99, 104897,  (2023)\\


 R. Baggio, O.U. Salman, L. Truskinovsky, "\textit{Inelastic rotations and pseudoturbulent plastic avalanches in crystals
}", Phys. Rev. E 107, 025004 – Published 24 February (2023)\\


C. Baruffi, A. Finel, Y. Le Bouar, B. Bacroix, O.U. Salman, "\textit{Atomistic simulations of temperature-driven microstructure formation in pure Titanium
}", Computational Materials Science 203, 111057, Editor's choice, (2022)\\

O.U. Salman, I.R. Ioanescu, "\textit{Tempering the mechanical response of FCC micro-pillars: an Eulerian plasticity approach
}",  Mechanics Research Communications, Volume 114, June , (2021)\\

O.U. Salman, L. Truskinovsky, "\textit{Delocalizing fracture
}", Journal of the Mechanics and Physics of Solids 154,  104517 , (2021)\\

OU Salman, R Baggio, B Bacroix, G Zanzotto, N Gorbushin, "\textit{Discontinuous yielding of pristine micro-crystals
}", Comptes Rendus. Physique 22 (S3), 1-48 , (2021)\\


P. Zhang, O.U. Salman, J. Weiss, L. Truskinovsky, "\textit{Fluctuations in crystalline plasticity
}", Comptes Rendus. Physique 22 (S3), 1-3707, (2021)\\


P. Zhang, O.U. Salman, J. Weiss, L. Truskinovsky, "\textit{Variety of scaling behaviors in nanocrystalline plasticity
}", Phys. Rev. E, 102, 023006, arXiv:2004.08579, hal-02644607, (2020)\\

Edited by I.R. Ionescu, S. Queyreau, C. R. Picu, O.U. Salman, "\textit{Mechanics and Physics of Solids at Micro- and Nano-Scales
}", ISBN: 978-1-786-30531-2 February (2020)\\

O.U. Salman, R. Baggio, "\textit{Homogeneous Dislocation Nucleation in Landau Theory of Crystal Plasticity}", appeared in the book \textit{Mechanics and Physics of Solids at Micro- and Nano-Scales
}, https://doi.org/10.1002/9781119687566.ch1, (2019)\\


R. Baggio, E. Arbib, P. Biscari, S. Conti, L. Truskinovsky, G. Zanzotto, and O.U. Salman, "\textit{Landau-Type Theory of Planar Crystal Plasticity
}", Phys. Rev. Lett. 123, 205501,  (2019)\\

O.U. Salman, G. Vitale, L. Truskinovsky, “\textit{Continuum theory of bending-to-stretching transition}“, Phys. Rev. E 100, 051001(R), Rapid Communication (2019)\\

O.U. Salman, B. Muite, A. Finel, “\textit{Origin of stabilisation of macrotwin boundaries in martensites}”, The European Physical Journal B 92 (1), 20, (2019) \\

C. Baruffi, A. Finel, Y. Le Bouar, B. Bacroix, O.U. Salman,  “\textit{Overdamped Langevin Dynamics simulations of grain boundary motion}”, Materials Theory, 3,4 (2019)\\

P. Franciosi, M.Spagnuolo, O.U. Salman, "\textit{Mean Green operators of deformable fiber networks embedded in a compliant matrix and property estimates}
", Continuum Mechanics and Thermodynamics, Volume 31, Issue 1, pp 101–132, (2019) \\

A. Ask, S. Forest, B. Appolaire, K. Ammar, O.U. Salman, "\textit{A Cosserat crystal plasticity and phase field theory for grain boundary migration}", Journal of the Mechanics and Physics of Solids 115, 167-194, (2018)\\

P. Zhang, O.U. Salman, J.-Y. Zhang, G. Liu, J. Weiss, L. Truskinovsky and J. Sun, “\textit{Taming intermittent plasticity at small scales}” Acta Mater, 128, 351-364, (2018)\\

M. Jedrychowski, B. Bacroix, O.U. Salman, J. Tarasiuk, S. Wronski, “\textit{Investigation of SIBM driven recrystallization in alpha Zirconium based on EBSD data and Monte Carlo modeling}“,  Materials Science and Engineering 89, 012029,  (2016)\\

A. Taloni, O.U. Salman, S. Zapperi, “\textit{Volume Changes During Active Shape Fluctuations in Cells}“,  Phys. Rev. Lett. 114, 208101, (2015) - (Highlighted in Nature Physics and APS Physics Journals)\\


O. Schylgo, O.U. Salman, A. Finel, “\textit{Martensitic phase transformations in Ni-Ti base shape memory alloys: Landau Theory}“, Acta Materialia 60, 6784-6792, (2012)\\


O.U. Salman, A. Finel, R. Delville, D. Schryvers, “\textit{The role of phase compatibility in martensite}”, J.Appl.Phys.111,103517, (2012)\\

O.U. Salman, L. Truskinovsky, “\textit{On the critical nature of the plastic flow: one and two dimensional models}”, International Journal of Engineering Science, 59:219.254, (2012)\\

O.U. Salman, L. Truskinovsky “\textit{Minimal discrete automaton behind crystal plasticity}“,  Phys. Rev. Lett. 106, 175503, (2011)\\

A. Finel, O.U. Salman, “\textit{Phase Field Methods: microstructures, mechanical properties and complexity}”, Comptes Rendus Physique, Vol. 11, Issues 3-4, p. 245-256, (2010) \\

 O.U. Salman, “\textit{Criticality in martensite}”, Mathematisches Forschungsinstitut Oberwolfach, DOI: 10.4171/OWR/2009/57 (2009) \\

R. Delville R. D. James, O.U. Salman, A. Finel and D. Schryvers“\textit{Transmission  electron  microscopy  study  of  low-hysteresis  shape  memory alloys }”, DOI:10.1051/esomat/200902005, ESOMAT 2009, 02005 (2009) \\

B. Muite, O.U. Salman, "\textit{Computations of geometrically linear and nonlinear Ginzburg-Landau models for martensitic pattern formation}", in ESOMAT 2009 edited by E. Sciences, P. Sittner, V. Paidar, L. Heller, H. Seiner, 587
p. 03008, (2009) \\

S. Ugur, O.U. Salman, G. Tepehan, F. Tepehan, O. Pekcan, \textit{Fluorescence study on $Al_2O_3$ polystyrene latex composite film formation},
Polymer composites 26 (3), 352-360, (2005)

%\NewPart{Publications in preparation}{}
%\hspace{3mm}
%\textbf{N. Gorbushin, O.U. Salman, L. Truskinovsky, "\textit{Marginal stability in Landau theory of crystal plasticity
%}" , (2020)}\\
%
%%\textbf{R. Baggio, O.U. Salman, L. Truskinovsky, B. Bacroix, "\textit{Symmetry-induced dislocation nucleation patterns 
%%}" , (2020)}\\
%
%\textbf{M .Degetier , O.U. Salman, D. Faurie "\textit{Substrate stifness dependance of fracture patterns in thin films  
%}" , (2020)}\\

%\textbf{E. Arbib, O.U. Salman, "\textit{A survey of Finite Element methods for crystal plasticity 
%}" , (2020)}\\


\NewPart{Conferences}{}
\hspace{3mm}

International Conference on processing and manufacturing  of advanced materials, Invited speaker, July 2 - 7, 2023
Vienna, Austria\\

"Plastic Turbulence", Invited speaker, Université Pierre et Marie Curie, Institut d'Alembert,  18 May (2022)\\

"Mécanique des Matériaux Architecturés", Invited speaker, Aussois, 17-21 January (2022)\\

"Landau theory of crystal plasticity: Application to nano-pillar plasticity", Invited speaker, Lund University, Mechanics Department,  15 October (2021)\\

"Avalanche dynamics and precursors of catastrophic events"
 3-8th of February, Les Houches, France, (2019)\\

"11th US–France Mechanics and physics of solids ), Speaker/Organiser, Paris, France, (2018)\\

"Phase Transformations and Plasticity in Crystals:
Atomistic to Continuum Models", Invited Speaker, Politecnico di Milano (https://www.fisi.polimi.it/it/home/news/57090), Milano, September 3-4, (2018)\\
 
“Mesoscopic model of dislocational plasticity accounting for crystal symmetry “, ECCM - ECFD 2018 6th European Conference on Computational Mechanics (Solids, Structures and Coupled Problems), Glasgow, UK, (2018)\\

“Mesoscopic model of dislocational plasticity accounting for crystal symmetry", The 13th World Congress in Computational Mechanics, NY, USA, (2018)\\

“The Landau theory of crystal plasticity  “, Mechanics and physics of solids at micro- and nano-scales, 11th US-France Symposium, Paris,  (2018)\\

%“Taming intermittent crystal plasticity“, LMS, Ecole Polytechnique, France, (2017)\\
%
%
%“Eulerian plasticity modelling of plasticity in micro-pillars“, COMPLAS, Barcelona, SPAIN, (2017)\\
%
%“Reversible and reconstructive martensitic phase transformations“, Réunion de la fédération mécanique francilienne, Fontainebleau, FRANCE, (2016)\\
%
%“Volume Changes During Active Shape Fluctuations in Cells“, School in Advances in Complex Systems, Como, ITALY, (2015)\\
%
%“Comparison of geometrically nonlinear and linear models for martensitic phase transformation”, ECCM, Paris, FRANCE, (2009)\\
%
%“Comparison of geometrically nonlinear and linear models for martensitic phase transformation”, ICOMAT- 08, Santa Fe, USA, (2008)\\
%
%“Avalanches in fluctuationless martensitic phase transformation”, ICOMAT-08, Santa Fe, USA, (2008)\\
%
%“Phase Field Modeling: Precipitation and Dynamics of the Martensitic Transition”, Annual TME Meeting, Orlando, USA, (2007)\\
%
%“Recent Advances in Phase Field Modeling, International Conferencen Advanced Materials” (ICAM 2007), Bangalore, INDIA, (2007)\\
%
%“Dynamics of Fluctuationless Martensitic Phase Transitions”, Euromat 2007, Nurenberg, GERMANY, (2007)\\
%
%“Theoretical and Numerical Comparison of Continuum Models for Martensitic Phase Transitions: Phase Field vs. Lagrangian Dynamics”, Phase-Field Models for the Evolution of Complex Structures - Institut Henry Poincarré, Paris, FRANCE, (2007)\\
%
%“Theoretical and Numerical Comparison of Continuum Models for Martensitic Phase Transitions: Phase Field vs. Lagrangian Dynamics”, Sixth Multimat Meeting, Prague, CZECH REPUBLIC, (2007)\\
%\hspace{3mm}
%
%“Phase Field Study of Martensitic Phase Transitions”, SPF Journées de la Matière Condensée (JMC10) Symposium: Modeling of defects in Materials, Toulouse, FRANCE, (2006)\\
%
%“Comparisons of a Geometrically Linear and a Geometrically Nonlinear Viscoelastic 2D Vectorial Model for Martensitic Pattern Formation”, Seventh Multimat Meeting, University of Oxford, ENGLAND, (2007)\\

\NewPart{Highlights}{}
\hspace{3mm}
The paper \textit{Landau-Type Theory of Planar Crystal Plasticity} published in the journal Physical Review Letters based on the Ph.D. work of R. Baggio that I co-supervised   is highlighted in the journal cover. The same work is also highlighted on the website of INSIS:

https://insis.cnrs.fr/fr/cnrsinfo/un-nouveau-modele-pour-etudier-la-plasticite-des-cristaux \begin{center}
   \includegraphics[scale=0.15]{IMG/cv123020.png}
 \end{center}
\par
The paper \textit{De-localisation of brittle fracture}, published in JMPS  is highlighted on the website of INSIS:

https://www.insis.cnrs.fr/fr/cnrsinfo/eviter-la-rupture-dun-materiau-en-redistribuant-la-deformation\\


The paper \textit{Cosserat crystal plasticity and phase field theory for grain boundary migration}, published in JMPS based on the post-doc work of A. Ask that I co-supervised is highlighted on the website of INSIS:

http://www.cnrs.fr/insis/recherche/actualites/2018/09/joints-grains-metal.htm\\

Received funding from Lundi University  Mechanics Department for a stay of two months to work with A. Ahadi, (2021).\\

Received funding from Bonn University for a stay of two weeks to work with S. Conti and M. Mungan, (2019).

\NewPart{Teaching}{}
\hspace{3mm}
 \textbf{Visiting Professor} – Lund University Mechanical Department, Sweden, (2022)\\

 \textbf{Invited Professor} – Romanian-French Summer School in Applied Mathematics, Sinaia, ROMANIA, (2018)\\

 \textbf{Lecturer} - Modelling of Functional Materials Master 2 Course (21H): Statistical mechanics, Theory of elasticity and Finite Element theory, ABAQUS Modelling Master 2 Course (21H), Université Sorbonne Paris Nord, FRANCE, (2015 - Present)\\
% \textbf{Teaching Assistant} - Finite Elements Course Engineering Department, Harvard University, USA, (2012) 
%\NewPart{Funding}{}
%\hspace{3mm}
\NewPart{Funding}{}
\hspace{3mm}
 \textbf{ANR: MESOCRYSP} - \textit{Discrete Mesoscale Plasticity}, 450000\euro~ for LSPM, \textbf{Project Partner, Leader of LSPM team, collaboration PMMH-Centre des Matériaux}, (11/2021-09/2025)\\


 \textbf{ANR-JCJC: ALIS} - \textit{Modeling of reconstructive phase transformations through Atomistically informed Landau theory with Infinite Symmetry}, 160000\euro~ for LSPM, \textbf{Project leader}, (03/2020-03/2023)\\
 
 \textbf{ANR-FWF: NANOFILM} - \textit{NANOARCHITECTED FILMS FOR UNBREAKABLE FLEXIBLE ELECTRONICS} (NANOFILM), 572K\euro~ in total, 142K\euro~ for  LSPM, Collaboration with PPRIME-CNRS, ESI Leoben (Austria), I will lead the phase-field modelling of fracture of thin films and supervise one post-doc, (03/2021-03/2024)\\
 
 \textbf{ANR-PRCI-SUMMIT} - Collaboration with Xian Jiaotong University, (29 500\euro~ for LSPM (265500\euro~  total)
 Euro), Project partner, (01/2018-12/2021)\\

 \textbf{Labex-Projet émergent} - Collaboration with M. Durant of Paris 7 University, UNIfied multi-scale framework for PLASticity of metals, amorphous and cellular materials, Project leader, (10/2018-10/2019)\\
 \textbf{ANR-I2T2M}, Project partner, (3K\euro~  total (01/2019-12/2022), \textbf{PICS-Pologne} - International Scientific Collaboration Project, CNRS (15K\euro~), Project manager, (2017 – 2020), \textbf{Visiting Professor }- Funding Paris 13 University (2,5K Euro), (2017), \textbf{Concours coup de pouce}- F2M – French Mechanics Federation (50K \euro~), Project Leader, (2015 - 2016)


%\NewPart{Organisation}{}
%\hspace{3mm}
%Mechanics and physics of solids at micro- and nano-scales, 11th US-France Symposium, Paris, (http://icacm2018.eu/), \textbf{Co-organizer/Speaker},  19-21 June 2018
%
%
\NewPart{Supervisions}{}

\section*{Current}
\hspace{3mm}
Post-Doc: \textbf{K. Ghosh}, Modelling of coupling of phase transition with crystal plasticity in Zirconium and Titanium through Landau theory ,
funded by ANR-JCJC, (18 Months), 2022-2024,\\

Post-Doc: \textbf{M. Terzi},  A low dimensional representation of 6D metric space via machine learning for the development of a new 3D Landau theory of crystal plasticity,
funded by ANR-MESOCRYSP, (24 Months), 2021-2023,\\


Ph.D: \textbf{J. Smiri},  Eulerian finite element modelling of ductile fracture in micro-scale crystals,
funded by Funded through Paris-Nord Sorbonne Université Scholarship, co-encadrant (36 Months), 2021-2024,\\

Ph.D: \textbf{M. Aissaoui},  Acoustic emission study of intermittent plasticity: the role of disorder, crystal orientation and symmetry,
funded by  Paris-Nord Sorbonne Université Scholarship, co-encadrant (36 Months), 2021-2024,\\

Ph.D: \textbf{G. Tejedor},  Modélisation numérique de la plasticité mésoscopique des cristaux par l'élasticité non linéaire invariante GL(3, Z) 
, funded by  ANR-MESOCRYSP, co-encadrant (36 Months), 2022-2025

\section*{Completed}\hspace{1mm}
Post-Doc \textbf{M. Degeiter},  Modelling of Fracture in thin films through phase-field approach,
funded by Labex, 2019-2022,\\

Post-Doc \textbf{N. Gorbushin},  Landau-type modeling of marginal stability in plasticity, funded by ANR-PRCI-SUMMIT, 2019-2021 \\

Post-Doc \textbf{N. Perchikov},  Reduction to an automaton of geometrically nonlinear Landau theory of crystal plasticity, funded by ANR-PRCI-SUMMIT, 2019-2020 \\


Ph.D, \textbf{R. Baggio}, Multi-Scale Modelling of Plasticity,
 Funded through Paris-Nord Sorbonne Université Scholarship, co-encadrant,  defended in  December, 2019, (post-doc in INRIA-Bordeaux and Université de Corse)\\

Ph.D, \textbf{M. Spagnuolo}, Modélisation continue des structures pantographiques, 
European- Inspire Scholarship, EU, co-encadrant,  defended in September 2019\\

Post-Doc, \textbf {A. Ask}, Modélisation de Recristallisation
Concours coup de pouce, Fédération francilienne de mécanique, FRANCE, funded by Concours coup de pouce de F2M, (2015-2016), (currently hired as research engineer in ONERA) \\

Ph.D,  \textbf{C. Baruffi}, Modélisation des mécanismes de vieillissement et d’endommagement dans les alliages métalliques, ONERA Scholarship, FRANCE, defended in December 2018 (currently  post-doc in EPFL Switzerland,  \textbf{received  le prix Jacques Dalla Torre 2019 for her Ph.D. work})\\

Ph.D, \textbf{P. Zhang}, supervision of the last year of his Ph.D. work in Paris , collaboration China/France,  defended in  2018, (currently hired as assistant professor in Xi’an Jiaotong University)\\

Stage M2, \textbf{F. Mansouri}, PNUS/LSPM (2016), \textbf{O. Benamor},  PNUS/LSPM (2020) 

\NewPart{Reviewer Tasks}{}
\hspace{3mm}
I reviewed  15 articles between 2017-2022 submitted in journals: \textbf{Physical Review Letters, Physical Review B, Physical Review E, Physical Review X, Physical Review Materials, Computational Materials Science, Philosophical Magazine, Mechanics of Materials}.

 I joined the Editorial Board of Interdisciplinary Physics (specialty section of Frontiers in Physics) as a Review Editor since 2021. 

cThe C++ code implementing the crystal plasticity model based on GL(2,Z) invariance is available:

https://bitbucket.org/oguzumut-cnrs/cute-crystal/src

The C++ code implementing the phase field  model of fracture in thin films   is available:

https://bitbucket.org/oguzumut-cnrs/uspn-fracture-fourier/src
\NewPart{Resumé en français}{}
\hspace{3mm}

L'extrême miniaturisation de la technologie moderne exige une meilleure compréhension de la mécanique non conventionnelle, dominée par les fluctuations, des matériaux à des échelles micro et nanométriques. A ces échelles, la nature de la plasticité médiée par les dislocations change radicalement : les échantillons métalliques submicroniques présentent une limite d'élasticité élevée, cependant corrompue par des fluctuations de déformation intermittentes compromettant les processus de mise en forme et mettant en danger la stabilité structurelle. L' un des objectives  de ma recherche   est de mieux comprendre ces phénomènes en développant des modèles mathématiques en allant  au de là des  outils de la mécanique continue classique comme la  théorie de Landau en grande déformation qui peut représenter adéquatement les symétries particulières des réseaux cristallins ou le milieu de Cosserat. L'une de ma contribution majeure concerne la mise en évidence et la caractérisation de l’intermittence des événements résultant du comportement collectif des défauts (dislocations, joints de grains, macles). Il apporte ainsi une explication physique aux observations expérimentales les plus récentes réalisées sur les grands instruments (ESRF Grenoble). L’innovation de la contribution tient à la conjugaison d’approches discrètes et continues pour aborder la description des phénomènes critiques lors de la déformation de la matière. j'ai a pour cela développé une mécanique des milieux continus raffinée, combinée à la méthode des champs de phase.
Le caractère générique de ces avancées scientifiques se manifeste par les applications variées aux matériaux de l’ingénieur : métaux et alliages, matériaux composites (fibres longues, matrice polymère) matériaux amorphes et cellulaires.


%\NewPart{Conferences}{}
%\hspace{3mm}
%"Les Matériaux Architecturés", Invited speaker, Aussois, (2021)\\
%
%
%"Avalanche dynamics and precursors of catastrophic events"
% 3-8th of February, Les Houches, France, (2019)\\
%
%"11th US–France Mechanics and physics of solids at macro- and nano-scales” symposium ICACM (International Center for Applied Computational Mechanics), Speaker, Paris, France, (2018)\\
%
%"Phase Transformations and Plasticity in Crystals:
%Atomistic to Continuum Models", Invited Speaker, Politecnico di Milano (https://www.fisi.polimi.it/it/home/news/57090), Milano, September 3-4, (2018)\\
% 
%“Mesoscopic model of dislocational plasticity accounting for crystal symmetry “, ECCM - ECFD 2018 6th European Conference on Computational Mechanics (Solids, Structures and Coupled Problems), Glasgow, UK, (2018)\\
%
%“Mesoscopic model of dislocational plasticity accounting for crystal symmetry", The 13th World Congress in Computational Mechanics, NY, USA, (2018)\\
%
%“The Landau theory of crystal plasticity  “, Mechanics and physics of solids at micro- and nano-scales, 11th US-France Symposium, Paris,  (2018)\\
%
%“Eulerian plasticity modelling of plasticity in micro-pillars“, COMPLAS, Barcelona, SPAIN, (2017)\\
%
%“Reversible and reconstructive martensitic phase transformations“, Réunion de la fédération mécanique francilienne, Fontainebleau, FRANCE, (2016)\\
%
%“Volume Changes During Active Shape Fluctuations in Cells“, School in Advances in Complex Systems, Como, ITALY, (2015)\\
%
%“Comparison of geometrically nonlinear and linear models for martensitic phase transformation”, ECCM, Paris, FRANCE, (2009)\\
%
%“Comparison of geometrically nonlinear and linear models for martensitic phase transformation”, ICOMAT- 08, Santa Fe, USA, (2008)\\
%
%“Avalanches in fluctuationless martensitic phase transformation”, ICOMAT-08, Santa Fe, USA, (2008)\\
%
%“Phase Field Modeling: Precipitation and Dynamics of the Martensitic Transition”, Annual TME Meeting, Orlando, USA, (2007)\\
%
%“Recent Advances in Phase Field Modeling, International Conferencen Advanced Materials” (ICAM 2007), Bangalore, INDIA, (2007)\\
%
%“Dynamics of Fluctuationless Martensitic Phase Transitions”, Euromat 2007, Nurenberg, GERMANY, (2007)\\
%
%“Theoretical and Numerical Comparison of Continuum Models for Martensitic Phase Transitions: Phase Field vs. Lagrangian Dynamics”, Phase-Field Models for the Evolution of Complex Structures - Institut Henry Poincarré, Paris, FRANCE, (2007)\\
%
%“Theoretical and Numerical Comparison of Continuum Models for Martensitic Phase Transitions: Phase Field vs. Lagrangian Dynamics”, Sixth Multimat Meeting, Prague, CZECH REPUBLIC, (2007)\\
%\hspace{3mm}
%
%“Phase Field Study of Martensitic Phase Transitions”, SPF Journées de la Matière Condensée (JMC10) Symposium: Modeling of defects in Materials, Toulouse, FRANCE, (2006)\\
%
%“Comparisons of a Geometrically Linear and a Geometrically Nonlinear Viscoelastic 2D Vectorial Model for Martensitic Pattern Formation”, Seventh Multimat Meeting, University of Oxford, ENGLAND, (2007)\\
%\NewPart{Collaborations}{}
%\hspace{3mm}
%\section*{National}
%\textbf{Plasticity:} I collaborate with \textbf{L. Truskinovsky (ESPCI/CNRS)}, \textbf{J. Weiss (CNRS)}
%focused on the study of mechanical properties of submicron crystals. The common project is organized around the co-supervision of a doctoral student funded by Paris Nord Université Sorbonne and 3 post-doctoral researchers funded by an ANR grant working on the development of Landau type of mesoscopic models of plasticity and experiments on nano- and micro-pillars. I also work with Professor \textbf{I.R. Ioanescu (LSPM)} with whom we develop an Eulerian continuum model of plasticity focused on the mechanical behavior of pillars under compression.\\
%
%\textbf{Fracture:} I develop together  with \textbf{L. Truskinovsky (ESPCI/CNRS)} a new model describing fracture in meta-materials and also in cells. I also started to work with \textbf{F. Cazes (LSPM)} to print the architecture materials that can be linked back to  our theoretical work with superior fracture properties. \\
%
%\textbf{Atomistic modeling of grain boundary migration and phase transformations:}
%This project is based on theoretical  development  introduced during the Ph.D. studiy of C. Baruffi that I co-supervised together with \textbf{A. Finel(LEM/Université Paris-Saclay)}, \textbf{Y. Le Bouar (LEM/Université Paris-Saclay)} and \textbf{B. Bacroix (LSPM)}. We developed  a modeling approach describing  the evolution of particle positions with an overdamped stochastic dynamics that we applied grain boundary migration and solid-solid phase transformations in Titanium.
%\section*{International}
%I collaborate with \textbf{G. Zanzotto (Padova University)}, \textbf{S. Conti (Bonn University)} and \textbf{P. Biscari (Politecnico Milano)} on mathematical aspect of modeling of crystal plasticity. I  also recently started to work with \textbf{M. Mungan, Bonn University} to apply the network theory description developed for amorphous materials to crystal plasticity.
%
%
%
%\NewPart{Detailed Summary of scientific research}{}
%Generally, I have wide ranging interests in the emergence of complexity/criticality and pattern formation in non-equilibrium externally driven inhomogeneous systems with applications to diverse problems in materials science such as plasticity, solid-solid phase transformations and fracture.  Particularly, my research is focused on the mathematical modeling and make use of computational methods to gain insight into underlying physics of these non-equilibrium dynamical systems. During the period of 2017-2020, My main research effort was concentrated  on plasticity  together with some contribution in the modeling of phase transformations and fracture as planned in my CNRS research project.   Below, I will present a summary of scientific research that I carried out during the last 3 years.
%
%\section{Plasticity}
%Considerable research efforts have been recently focused on the study of mechanical properties of submicron crystals. It was found that the deformation mechanisms, which we habitually associate with dislocation plasticity, change dramatically once the sample size is reduced below the micrometer range. Strength of such crystals was shown to be size-dependent, with stress-strain response exhibiting pronounced intermittency and scale invariance over a wide range of scales, independently of crystal symmetry. Both measured and computed scaling exponents were shown to feature highly unusual size dependence.
%Moreover, even though plasticity at macroscale is generally associated with ductility, crystal plasticity at submicron scales exhibits major stress drops or strain bursts reminiscent of brittle fracture. Brittleness, usually attributed to dislocation-free crystals, reappears in nanoparticles and nanopillars that seem to be “breaking plastically” by generating a large number of globally correlated dislocations. The implied system-size events hinder our ability to control plastic deformation at submicron scale and compromise the reliable functioning of ultra-small machinery. 
%\begin{figure}[H]
%\begin{center}
%\includegraphics[scale=0.8]{fig11.png}
%\caption{SEM images of [112]-orientated micropillars after compression: (a) 500 nm, (b) 1000 nm, and (c) 1500 nm. The marked slip traces in panels (a) and (b) indicate the locally single-slip nature of the plastic flow in 500 nm and 1000 nm crystals.}
%\end{center}
%\label{p1}
%\end{figure}
%I collaborated with \textbf{J. Weiss (CNRS)  and L. Truskinovsky (CNRS/ESPCI)} on experimental and theoretical aspects of crystal plasticity. From experimental point of view, our goal was to propose new manufacturing protocols to avoid plastic localisation that we refer to as "breaking plastically". In our work published in Acta Mater, 128, 351-364, (2017),  we explored how the introduction of disorder, such as solutes or precipitates, used in classical metallurgy to harden materials, can be used at sub-micron scales to suppress intermittent fluctuations. Besides, we presented a mean-field type of model to present a theoretical framework to  conceptualize the observed behavior. Afterward, we received an ANR funding (ANR-PRCI-SUMMIT) and studied theoretically and experimentally size and disorder effects in BCC crystals, (see, Variety of scaling behaviors in nanocrystalline plasticity, Phys. Rev. E, 102, 023006 and Fig. \ref{p1} and Fig. \ref{p2}). Our main result is that the nonuniversality of submicron  plasticity and the inferred brittleness of ultrasmall crystals can be conceptualized as a multistage crossover from spin-glass marginality, characteristic of very small, almost defect-free crystals, to the criticality of larger crystals associated with a brittle-to-ductile (BD) transition.
%\begin{figure}[H]
%\includegraphics[scale=0.65]{fig12.png}
%\caption{After-yield spatial patterns with respect to the strength of the quenched disorder}
%\label{p2}
%\end{figure}
%Despite the fact that our simplified theoretical model captures the main trends observed in experiments, it falls short in reproducing special dislocation patterns. Therefore, it has to be extended  to the full tensorial theory based on nonlinear elasticity to  obtain a more detailed description of crystal plasticity. This  extension will allow one to distinguish between different crystallographic classes and different configurations of dislocation  cores. Such a model should be able to reproduce dislocation  walls and generate cell structures with realistic size distribution. A 3D theory of this type can also capture dislocation  cross-slip and address hardening behavior. To study surface  effects directly, one would need to come up with the boundary  conditions allowing for dislocation nucleation on the surface.  We recently sketched a path in this general direction  in our papers O.U. Salman, R. Baggio, "\textit{Homogeneous Dislocation Nucleation in Landau Theory of Crystal Plasticity}", appeared in the book \textit{Mechanics and Physics of Solids at Micro- and Nano-Scales}, (2019) and R. Baggio, E. Arbib, P. Biscari, S. Conti, L. Truskinovsky, G. Zanzotto, and O.U. Salman, "\textit{Landau-Type Theory of Planar Crystal Plasticity
%}", Phys. Rev. Lett. 123, 205501,  (2019). Moreover, we also formulated a new way of coupling plasticity with cracks in our meso-scale Landau-type model (R. Baggio, Ph.D. thesis). This new formulation resulted in a significant reduction of computational time and can render possible to theoretically understand the link between dislocation patterning and macroscopic fracture during cycling loading experiments, see Fig. \ref{p3}.  Our post-doc N. Gorbushin that I co-supervise with \textbf{J. Weiss and L. Truskonivsky} continues to work on this subject. He will be funded  until the end of 2021. I also continuously collaborate with  \textbf{S. Conti (Professor of Mathematics, University of Bonn)} and\textbf{ M. Mungan (Professor of Physics, University of Bonn)} and \textbf{I.R. Ioanescu (Professor in PNSU/LSPM}) on other aspects of crystal plasticity such as a network theory description of dislocation interactions and Eulerian description of plastic flow. 
%
%
%
%\begin{figure}[H]
%\begin{center}
%\includegraphics[scale=0.65]{fig14.png}
%\caption{: From top to bottom, post instability patterns   are shown in Landau theory and Molecular statics, respectively. Edge dislocations are observed in the first two cases, together with voids and other volumetric defects.}
%\label{p3}
%\end{center}
%\end{figure}
%
%\section{Fracture in meta-materials}
%Brittle materials fail macroscopically through the catastrophic growth of a system-size crack which usually originates from a microscopic flaw. The material away from the crack tip remains effectively unloaded and the whole breaking process is characterized by extreme strain localization. A prototype of brittle material is shown in Fig.\ref{chains}(a). Mathematically, such localization can be traced to the loss of ellipticity of continuum elasticity, see see Fig. \ref{localisation}. Physically, it is a consequence of material softening which, in turn, is a manifestation of the fundamental nonconvexity of the inter-particle interaction. Ultimate strain localization in brittle materials is the source of a host of undesirable issues from low fracture toughness  to dramatic dynamic effects associated with abrupt fragmentation.
%
%The main challenge in delocalizing brittle fracture is to avoid unstable damage propagation driven by stress concentration. The usual approach is to engineer a set of obstacles or dissipation centers promoting the uniformity of damage. The idea is that individual lock-up mechanisms would cause local saturation of damage and drive subsequent failure elsewhere. Usually cracks are channeled by such mechanisms into configurations where some structural conditions impede their unstable propagation and trap them into stable configurations. The geometry of the implied local reinforcements can be tailored to achieve high energy absorption so that the resulting nominally brittle materials can rival their ductile analogs.
%Many particular strategies for creating effective ductility have been proposed in the literature. If we exclude explicit ’plastification’, an incomplete list of nominally ’brittle’ approaches include toughening by micro-cracking, rigidity mitigation, transformational toughening  and the use of multi-level damage.
%
%In this work we proposed a principally different approach to fracture delocalization involving nonlocal interactions which diminish stress concentration by redistributing the load in a softening material. Such interactions can be transmitted by a system of effective backbones whose elasticity is bending dominated. We propose a particular mesoscopic design of such fracture delocalizing system which can be viewed as an architectured meta-material, see Fig. \ref{chains}(b). This design has a bio-mimetic nature and structurally imitates some known biological prototypes.
%
%\begin{figure}[H]
%\begin{center}
%\includegraphics[scale=0.5]{fig1.pdf}
%\caption{(a) The system consists of a 1D array of particles connected by Lennard-Jones  type springs, (b) Pantographic beam with breakable springs connecting the mid-points.}
%\label{chains}
%\end{center}
%\end{figure}
%
%More precisely, we show that the de-localization of damage can be facilitated if failure in the initially stretching dominated conventional elastic material is accompanied by the progressive switch to bending dominated nonlocal elasticity. More specifically, we develop a prototype model of a meta-material with nominally brittle constituents at the micro-scale which is engineered to transit from the conventional elastic behavior at small level of damage to bending dominated elastic response at large level of damage. We show that such transition ensures that instead of failure localization, anticipated in stretching dominated regime, the bending elasticity redistributes the stress globally so that every single brittle sub-element of such effective material-structure eventually breaks. The crucial role in making such variable-connectivity design possible is played by our ability to build a high contrast composite structure where failure of elements ensuring rigidity leads to the emergence of a floppy modes stabilized exclusively by bending, see Fig. \ref{delocalisation}. \textbf{All this has recently become possible due to the developments in additive manufacturing techniques}.
%
%We already used the ideas mentioned above to explain the transition from bending-dominated to stretching-dominated elastic response in semi-flexible fibrous networks that plays an important role in the mechanical behavior of cells and tissues, see Salman, Vitale, Truskinovsky, Continuum theory of bending-to-stretching transition}“, Phys. Rev. E 100, 051001(R), Rapid Communication (2019). I will continue to collaborate with \textbf{L. Truskinovsky (ESPCI-CNRS)} on the theoretical/mathematical aspects of fracture and extend our one dimensional theoretical framework to higher dimensions. On the other hand, we recently acquired a 3D printer in our lab LSPM using which we can print the architecture materials that can be linked back to  our theoretical work with superior fracture properties. I am collaborating with \textbf{Fabien Cazes (Maitre de Conference, LSPM/PNUS)} on this aspect.
%
%
%
%
%
%\begin{figure}[H]
%\begin{center}
%\includegraphics[scale=0.3]{fig3_neww.pdf}
%\caption{Crack localisation in a standard chain: (a) the equilibrium elastic energy $F(\bar\epsilon)$ and, (b) the equilibrium stress  $\sigma(\bar\epsilon) =\partial F(\epsilon)|_\epsilon$. The dashed black line corresponds to  the homogeneous solution and, (c) strain profiles for different values of loading parameter. Here $N=50$.}
%\label{localisation}
%\end{center}
%\end{figure}
%
%\begin{figure}[H]
%\begin{center}
%\includegraphics[scale=0.3]{figbendingg.pdf}
%\caption{Crack de-localisation in reinforced discrete chain: (a) the equilibrium elastic energy $F(\bar\epsilon)$ and, (b) the equilibrium stress  $\sigma(\bar\epsilon) =\partial F(\epsilon)|_\epsilon$. The dashed black line corresponds to  the homogenous solution and, (c)  strain profiles for different values of loading parameter. Here $N=50$.}
%\end{center}
%\label{delocalisation}
%\end{figure}
%
%\section{Fracture in thin films}
%Devices fabricated on flexible substrates such as polymers have been widely studied because
%of their remarkable potential for new applications requiring
%non-planar functional systems. During the last years, several electronic devices, such as flexible organic solar cells, light-emitting diodes and transistors, have been realized
%on top of various polymeric substrates and even on paper. In
%addition, the increasingly complex forms of consumer electronics and smart textiles will benefit from the magneto-electric features of these systems by their adaptability to non-planar surfaces (conformability). Most objects in our daily life are not flat and rigid, rather flexible and pliable, with the appropriate examples being sheets of paper and textiles. Integration of functional features into the flexible objects may provide numerous users with a number of benefits, and also gives rise to new applications such as flexible displays, rollable hard drives, and wearable sensors. The time is now ripe for  designing flexible systems with high potential for the future. Moreover, polymers are known to be lightweight and cheap. Hence, it is advantageous to use them as substrates for functional
%nanostructures; this requires adaption of the existing expertise, technologies and techniques to the
%polymer substrates.
%
%From a general point of view, the main limitation of these flexible/stretchable systems is their often too weak durability. If the polymer substrates are adapted to large strains, the inorganic films and
%interconnecting lines which carry the system functionality are intrinsically much more fragile. When a
%metal is deposited on a flexible substrate, it is subjected to high mechanical stresses due to the
%stretching or curvature of the entire system and the interfacial metal/polymer adhesion. The problem of early fracture of the thin films is not yet solved. Unfortunately, materials in the form of thin films are rarely ductile. The few exceptions are a few face centered cubic materials such as copper or gold. More recently, research on amorphous aluminum oxide has shown that this material can be very
%ductile in the form of a flawless thin film. However, these solutions based on the intrinsic structure
%of the material will not be able to scale all the functionalities required for flexible systems (optical,magnetic, and electrical). Thus, we believe that research on multi-materials with nanometric
%architecture will allow developing systems showing à la carte functionalities and support large strains.
%
%To this end, I started to collaborate with \textbf{D. Faurie (Professor, LSPM/PNUS)}, who has a great amount of insight into thin films' mechanical/fracture properties. We hired M. Degeiter as a post-doc in 2019 to start developing a phase-field model of fracture for thin-films on flexible substrates. As usual, we started with an analytically tractable 1D version of our full-model and performed a bifurcation analysis to predict the energy-minimizing crack nucleation patterns. Afterward,  we obtained global and local minimizers  by solving the set of mechanical equilibrium phase-field equations numerically, see  Fig.\ref{cracks}. This initial and mostly conceptual analytical/numerical work will guide our future work in higher dimensions. M. Degeiter will continue his post-doc through our recently funded ANR-FWF: NANOFILM  (I am co-partner and responsible for theoretical development within the project) starting from October 2020.
%\begin{figure}[H]
%\begin{center}
%\includegraphics[scale=0.35]{fig13.png}
%\caption{Global and local minimizers of a phase field model  for thin-films: (a) 
%energy difference between the actual configuration and the homogeneous configuration for all energy branches, (b-c) strain-stress relation, (d-i) strain and phase-field profiles for lowest energy-branch associated with nucleation of 3 cracks in thin-films on a rigid substrate. }
%\end{center}
%\label{cracks}
%\end{figure}
%
%\section{Atomistic modeling of grain-boundary migration and solid-solid phase transformations}
%We developed a modeling approach describing  the evolution of particle positions with an overdamped stochastic dynamics. In contrast to Newtonian dynamics, particle positions are treated as stochastic variables which follow a first-order in time dynamics that do not explicitly incorporate high-frequency vibrations of the crystalline grid (phonons), which typically limit to few nanoseconds the time scale of classical Molecular Dynamics. The chaotic nature of the Newtonian dynamics, which in the long time drives the system to a stochastic equilibrium state, is recovered in the first-order in time dynamics through the use of an additive noise term, carefully chosen to guarantee that the system converges to the correct thermodynamical state in the long-time limit. We applied this method to the to study the shrinking of a circular grain embedded inside a crystal \textit{Overdamped Langevin Dynamics simulations of grain boundary motion}”, Materials Theory, 3,4 (2019) and BCC-to-HCP phase phase transformation observed in Titanium, see Fig. \ref{ptt} .  The collaboration is structured by the co-supervision of a doctoral student (C. Baruffi) funded by ONERA/Université Sorbonne Paris Nord. The experimental studied on both grain boundary  migration and phase transformations has been studied in LSPM, which are now supported by modeling approaches.
%\begin{figure}[H]
%\begin{center}
%\includegraphics[scale=0.4]{graphical_abstract.pdf}
%\end{center}
%\caption{Microstructure formation observed during temeprature-driven BCC-HCP transition and the evolution of volume fractions. We identified the anti-phase boundary formation or self-accommodation patterns in Titanium with respect to applied external mechanical constraints.   }
%\label{ptt}
%\end{figure}
%%% References
%%% ------------------------------------------------------------

\end{document}
